% Κοινή κατασκευή AttentionBox / ThinkBox / HistoryBox.
% Οι δηλώσεις στο τέλος αλλάζουν μόνο φόντο, τόνο έμφασης, εικονίδιο,
% και ετικέτα. Γεωμετρία, τίτλος, πλευρική γραμμή και συμπεριφορά
% αλλαγής σελίδας υλοποιούνται μία φορά για όλα τα κουτιά.
% Προαπαιτούμενα: BookCallout, BookCalloutCenteredIcon, BookMediumLabelFont
% και οι τρεις σχεδιαστικές εντολές εικονιδίων του universalbook.sty.

\providecommand{\BookCalloutLabelSize}{9.2}
\providecommand{\BookCalloutLabelLeading}{11.5}

% Ο τίτλος είναι σώμα του κουτιού (όχι δεύτερο ανεξάρτητο title box).
% Το προαιρετικό θέμα διατηρεί τη σύνταξη HistoryBox[Θέμα].
\newcommand{\BookCalloutHeading}[3]{%
  \noindent{\BookMediumLabelFont
    \fontsize{\BookCalloutLabelSize}{\BookCalloutLabelLeading}\selectfont
    \color{#1}#2\ifstrempty{#3}{}{ — #3}}\par\smallskip\small
}

% Πέντε παράμετροι: φόντο / έμφαση / εικονίδιο / τίτλος / θέμα.
% Το ίδιο overlay χρησιμοποιείται σε ενιαίο και πρώτο σπασμένο τμήμα.
% Τα επόμενα τμήματα κρατούν τη γεωμετρία χωρίς να επαναλαμβάνουν icon.
\tcbset{BookCalloutRender/.style n args={5}{%
  BookCallout,
  colback=#1,
  borderline west={1.35pt}{0pt}{#2},
  before upper={\BookCalloutHeading{#2}{#4}{#5}},
  overlay unbroken={\BookCalloutCenteredIcon{#3}},
  overlay first={\BookCalloutCenteredIcon{#3}}
}}

% Λεπτά περιβλήματα με την παλιά ακριβώς δημόσια σύνταξη.
% Η μόνη διαφορά του δεύτερου είναι το προαιρετικό θέμα.
\newcommand{\BookDeclareCallout}[5]{%
  \newtcolorbox{#1}{BookCalloutRender={#2}{#3}{#4}{#5}{}}%
}
\newcommand{\BookDeclareTitledCallout}[5]{%
  \newtcolorbox{#1}[1][]{BookCalloutRender={#2}{#3}{#4}{#5}{##1}}%
}

\BookDeclareCallout{AttentionBox}{BookAlertBackground}{BookAlert}{\BookCalloutWarning}{\LabelAlert}
\BookDeclareCallout{ThinkBox}{BookThinkBackground}{BookThinkAccent}{\BookCalloutBulb}{\LabelThink}
\BookDeclareTitledCallout{HistoryBox}{BookHistorySoft}{BookHistory}{\BookCalloutOpenBook}{\LabelHistory}
